\lecture{22}{Дополненная функция Лагранжа}

\sourcevideo
    {https://www.youtube.com/playlist?list=PLOzRYVm0a65fhKDS8cYIC7YCuVGJ4FOJx}
    {Week 6: Lecture 22: Augmented Lagrangian}

Рассмотрим задачу с ограничениями-равенствами
\[
    \min_{x\in\mathbb R^n} f(x)
    \qquad
    \text{при условии}
    \qquad
    h(x)=0,
\]
где
\[
    h\colon\mathbb R^n\to\mathbb R^m.
\]
Обычная функция Лагранжа равна
\[
    \mathcal L(x,\nu)
    =
    f(x)+\nu^{\mathsf T}h(x),
\]
а дополненная функция Лагранжа определяется формулой
\[
    \mathcal L_c(x,\nu)
    =
    f(x)
    +
    \nu^{\mathsf T}h(x)
    +
    \frac c2\lVert h(x)\rVert^2,
    \qquad
    c>0.
\]
Квадратичный член не изменяет значение функции на допустимом
множестве, но добавляет кривизну в направлениях, нарушающих
ограничение. Эта конструкция лежит в основе метода множителей и ADMM.

\section{Штрафная и дополненная постановки}

Следует различать исходную задачу
\[
    \min_x f(x)
    \qquad
    \text{при условии}
    \qquad
    h(x)=0,
\]
эквивалентную задачу
\[
    \min_x
    \left\{
        f(x)+\frac c2\lVert h(x)\rVert^2
    \right\}
    \qquad
    \text{при условии}
    \qquad
    h(x)=0
\]
и безусловную штрафную задачу
\[
    \min_x
    \left\{
        f(x)+\frac c2\lVert h(x)\rVert^2
    \right\}.
\]
Первые две задачи имеют одинаковые допустимые точки и одинаковые
значения целевой функции на них. В третьей задаче ограничение удалено,
поэтому при конечном \(c\) её решение в общем случае не обязано быть
допустимым. При подходящих предпосылках выполняется лишь
\[
    h(x_c)\to0
    \qquad
    \text{при }c\to\infty.
\]

Дополненная функция Лагранжа сочетает квадратичный штраф с линейным
членом \(\nu^{\mathsf T}h(x)\). Именно обновление множителя позволяет
методу множителей работать при конечном и умеренном штрафном
параметре.

Для линейного ограничения
\[
    h(x)=Ax-b,
    \qquad
    A\in\mathbb R^{m\times n},
\]
штраф имеет градиент и Гессиан
\[
    \nabla_x\frac c2\lVert Ax-b\rVert^2
    =
    cA^{\mathsf T}(Ax-b),
\]
\[
    \nabla_x^2\frac c2\lVert Ax-b\rVert^2
    =
    cA^{\mathsf T}A.
\]
Матрица \(A^{\mathsf T}A\) положительно полуопределена, поскольку
\[
    y^{\mathsf T}A^{\mathsf T}Ay
    =
    \lVert Ay\rVert^2
    \ge0.
\]
Она положительно определена тогда и только тогда, когда
\[
    \ker A=\{0\},
\]
то есть \(A\) имеет полный столбцовый ранг. В этом случае штраф
сильно выпуклый с константой
\[
    c\sigma_{\min}(A)^2.
\]
Если \(\ker A\ne\{0\}\), штраф не создаёт кривизну вдоль ядра.
Для положительной определённости \(A^{\mathsf T}A\) необходим полный столбцовый ранг \(A\).

\section{Роль множителя и квадратичный пример}

На допустимом множестве
\[
    \mathcal L_c(x,\nu)
    =
    \mathcal L(x,\nu)
    =
    f(x).
\]
При фиксированном произвольном \(\nu\) минимизаторы
\(\mathcal L(\cdot,\nu)\) и \(\mathcal L_c(\cdot,\nu)\) могут
различаться и не обязаны совпадать с решением исходной задачи.
Совпадение обеспечивается при правильном множителе и соответствующих
условиях оптимальности.

Рассмотрим задачу
\[
    \min_{x_1,x_2}
    \frac12(x_1^2+x_2^2)
    \qquad
    \text{при условии}
    \qquad
    x_1=1.
\]
Её решение равно
\[
    x^\star=
    \begin{pmatrix}
        1\\
        0
    \end{pmatrix}.
\]
Функция Лагранжа имеет вид
\[
    \mathcal L(x_1,x_2,\nu)
    =
    \frac12(x_1^2+x_2^2)+\nu(x_1-1).
\]
Из условий стационарности и допустимости
\[
    x_1+\nu=0,
    \qquad
    x_2=0,
    \qquad
    x_1=1
\]
получаем
\[
    \nu^\star=-1.
\]
Дополненная функция равна
\[
\begin{aligned}
    \mathcal L_c(x_1,x_2,\nu)
    ={}&
    \frac12(x_1^2+x_2^2)
    +\nu(x_1-1)
    \\
    &+
    \frac c2(x_1-1)^2.
\end{aligned}
\]
При фиксированном \(\nu\) её минимизатор удовлетворяет
\[
    x_2=0,
    \qquad
    x_1+\nu+c(x_1-1)=0,
\]
то есть
\[
    x_1(c,\nu)
    =
    \frac{c-\nu}{c+1}.
\]
При \(\nu=\nu^\star=-1\) получаем \(x_1=1\) для любого \(c\ge0\),
тогда как при произвольном фиксированном множителе
\[
    x_1(c,\nu)\to1
    \qquad
    \text{при }c\to\infty.
\]
Тем самым к решению можно приближаться либо увеличением штрафа, либо
уточнением двойственной переменной. Метод множителей использует второй
механизм.

\section{KKT-условия и локальная кривизна}

Пусть \(x^\star\) является локальным решением, а якобиан ограничения
равен
\[
    J_h(x)
    =
    \begin{pmatrix}
        \nabla h_1(x)^{\mathsf T}\\
        \vdots\\
        \nabla h_m(x)^{\mathsf T}
    \end{pmatrix}.
\]
При условии регулярности существует множитель \(\nu^\star\), для
которого выполняются KKT-условия
\[
    h(x^\star)=0,
\]
\[
    \nabla f(x^\star)
    +
    J_h(x^\star)^{\mathsf T}\nu^\star
    =0.
\]
Эквивалентно,
\[
    \nabla_x\mathcal L(x^\star,\nu^\star)=0.
\]
Поскольку \(h(x^\star)=0\), добавление штрафа сохраняет
стационарность:
\[
    \nabla_x\mathcal L_c(x^\star,\nu^\star)=0.
\]

Для произвольной точки
\[
\begin{aligned}
    \nabla_{xx}^2\mathcal L_c(x,\nu)
    ={}&
    \nabla_{xx}^2\mathcal L(x,\nu)
    +
    cJ_h(x)^{\mathsf T}J_h(x)
    \\
    &+
    c\sum_{i=1}^{m}h_i(x)\nabla^2h_i(x).
\end{aligned}
\]
В допустимой точке последний член исчезает:
\[
\begin{aligned}
    \nabla_{xx}^2
    \mathcal L_c(x^\star,\nu^\star)
    ={}&
    \nabla_{xx}^2
    \mathcal L(x^\star,\nu^\star)
    \\
    &+
    cJ_h(x^\star)^{\mathsf T}J_h(x^\star).
\end{aligned}
\]

Касательное пространство линеаризованного ограничения равно
\[
    \mathcal T(x^\star)
    =
    \ker J_h(x^\star).
\]
Если \(d\in\mathcal T(x^\star)\), то
\[
    d^{\mathsf T}
    J_h(x^\star)^{\mathsf T}J_h(x^\star)d
    =
    \lVert J_h(x^\star)d\rVert^2
    =
    0.
\]
Следовательно, квадратичный штраф не добавляет кривизны в касательных
направлениях. Там положительность должна обеспечиваться исходной
функцией Лагранжа. В нормальных направлениях штраф добавляет
\[
    c\lVert J_h(x^\star)d\rVert^2.
\]

\begin{definition}[Достаточное условие второго порядка]
Пара \((x^\star,\nu^\star)\) удовлетворяет достаточному условию
второго порядка, если выполняются KKT-условия и
\[
    d^{\mathsf T}
    \nabla_{xx}^2\mathcal L(x^\star,\nu^\star)d
    >0
\]
для каждого ненулевого
\[
    d\in\ker J_h(x^\star).
\]
\end{definition}

\section{Локальная точность дополненной функции}

\begin{lemma}
Пусть \(P\) и \(Q\) симметричны, \(Q\succeq0\), а
\[
    y^{\mathsf T}Py>0
\]
для каждого ненулевого \(y\in\ker Q\). Тогда существует
\(\overline c\ge0\), такое что
\[
    P+cQ\succ0
\]
для всех \(c\ge\overline c\).
\end{lemma}

Матрица \(P\) уже положительна в направлениях, где \(Q\) обращается
в нуль. В остальных направлениях достаточно большой положительный
член \(cQ\) компенсирует возможную отрицательную кривизну \(P\). В
рассматриваемой задаче
\[
    P=
    \nabla_{xx}^2\mathcal L(x^\star,\nu^\star),
    \qquad
    Q=
    J_h(x^\star)^{\mathsf T}J_h(x^\star),
\]
причём
\[
    \ker Q=\ker J_h(x^\star).
\]

\begin{theorem}[Локальная точность]
Пусть \(f\) и \(h\) дважды непрерывно дифференцируемы, а пара
\((x^\star,\nu^\star)\) удовлетворяет KKT-условиям и достаточному
условию второго порядка. Тогда существует \(\overline c\ge0\), такое
что для каждого \(c\ge\overline c\) точка \(x^\star\) является
строгим локальным минимумом функции
\[
    x\longmapsto\mathcal L_c(x,\nu^\star).
\]
\end{theorem}

\begin{proof}
KKT-условия и допустимость дают
\[
    \nabla_x\mathcal L_c(x^\star,\nu^\star)=0.
\]
В этой точке
\[
\begin{aligned}
    \nabla_{xx}^2
    \mathcal L_c(x^\star,\nu^\star)
    ={}&
    \nabla_{xx}^2
    \mathcal L(x^\star,\nu^\star)
    \\
    &+
    cJ_h(x^\star)^{\mathsf T}J_h(x^\star).
\end{aligned}
\]
По достаточному условию второго порядка первая матрица положительна на
ядре второй. Лемма обеспечивает положительную определённость полного
Гессиана при всех достаточно больших \(c\). Поэтому \(x^\star\) —
строгий локальный минимум дополненной функции при фиксированном
множителе \(\nu^\star\).
\end{proof}

Теорема локальна. Она не утверждает, что для произвольного \(\nu\)
глобальный минимум \(\mathcal L_c(\cdot,\nu)\) совпадает с
\(x^\star\).

В выпуклой задаче с аффинным ограничением результат усиливается. Если
\((x^\star,\nu^\star)\) является седловой точкой обычной функции
Лагранжа, то
\[
    \mathcal L(x^\star,\nu)
    \le
    \mathcal L(x^\star,\nu^\star)
    \le
    \mathcal L(x,\nu^\star).
\]
Поскольку штраф неотрицателен и равен нулю в \(x^\star\),
\[
    \mathcal L_c(x^\star,\nu)
    \le
    \mathcal L_c(x^\star,\nu^\star)
    \le
    \mathcal L_c(x,\nu^\star).
\]
Следовательно, та же пара остаётся седловой точкой дополненной
функции Лагранжа.

\section{Метод штрафов и метод множителей}

Метод квадратичных штрафов решает последовательность задач
\[
    x^{k+1}
    \in
    \operatorname*{argmin}_x
    \left\{
        f(x)
        +
        \frac{c_k}{2}\lVert h(x)\rVert^2
    \right\},
    \qquad
    c_k\to\infty.
\]
При подходящих предпосылках \(h(x^k)\to0\), но рост \(c_k\) ухудшает
обусловленность. Для линейного ограничения Гессиан содержит
\(c_kA^{\mathsf T}A\): кривизна в нормальных направлениях быстро
растёт, а в касательных может почти не изменяться.

Метод множителей использует итерации
\[
    x^{k+1}
    \in
    \operatorname*{argmin}_x
    \mathcal L_c(x,\nu^k),
\]
\[
    \nu^{k+1}
    =
    \nu^k+c\,h(x^{k+1}).
\]
Второе равенство корректирует множитель пропорционально невязке.
Для \(h(x)=Ax-b\)
\[
    \nu^{k+1}
    =
    \nu^k+c(Ax^{k+1}-b).
\]
Двойственная переменная аккумулирует нарушение допустимости и меняет
следующую прямую подзадачу. В отличие от чистого метода штрафов,
параметр \(c\) часто можно оставить фиксированным. Постоянное увеличение \(c\) возможно, но не является обязательным свойством метода множителей.

Слишком большой штраф нежелателен. Гессиан прямой подзадачи содержит
\[
    cJ_h(x)^{\mathsf T}J_h(x),
\]
поэтому растут число обусловленности и липшицева константа градиента,
уменьшается допустимый шаг явных методов и усиливается чувствительность
к ошибкам вычисления невязки. На практике выбирают умеренный \(c\) и
корректируют его по прогрессу допустимости.

Логарифмический барьер для неравенства \(g_i(x)\le0\) имеет вид
\[
    -\tau\sum_i\log\bigl(-g_i(x)\bigr).
\]
Он определён только при строгой допустимости и не позволяет пересекать
границу. Квадратичный штраф определён по обе стороны равенства и
допускает недопустимые промежуточные точки. Для равенства
\(h(x)=0\) логарифмический барьер непосредственно неприменим, поскольку
допустимое множество не имеет внутренности.

\section{Границы модели и переход к ADMM}

Содержательная идея состоит в следующем: достаточно большой
штраф компенсирует отрицательную кривизну в направлениях, меняющих
ограничение, но не влияет на касательное пространство.

При анализе дополненной функции сначала проверяются KKT-условия, затем
определяется
\[
    \mathcal T(x^\star)=\ker J_h(x^\star),
\]
и на нём проверяется условие второго порядка. После этого формула
\[
\begin{aligned}
    \nabla_{xx}^2
    \mathcal L_c(x^\star,\nu^\star)
    ={}&
    \nabla_{xx}^2
    \mathcal L(x^\star,\nu^\star)
    \\
    &+
    cJ_h(x^\star)^{\mathsf T}J_h(x^\star)
\end{aligned}
\]
позволяет выбрать достаточно большое, но численно приемлемое значение
\(c\).

Если прямая переменная разбита на несколько блоков и минимизация
дополненной функции выполняется по ним последовательно, возникает
метод чередующихся направлений множителей, или ADMM. Блочная структура
разделяет исходную задачу на подзадачи и далее используется для
построения распределённых алгоритмов.
